% © 2025 Ing. David Jaroš — CC BY-NC-ND 4.0
%
% This work is licensed under a Creative Commons Attribution-NonCommercial-NoDerivatives
% 4.0 International License (CC BY-NC-ND 4.0).
%
% License History: Earlier drafts (up to v0.3) were released under CC BY 4.0.
% From v0.4 onward, all material is released under CC BY-NC-ND 4.0 to protect
% the integrity of the theoretical work during ongoing academic development.
%
% Three Generations Research Track — ST-2: Obstruction Theorem
% Status: Complete proof
% Version: 1.0
% Date: 2026-03-03

\documentclass[12pt,a4paper]{article}

\usepackage{amsmath,amssymb,amsthm}
\usepackage{hyperref}
\usepackage{geometry}
\usepackage{enumitem}
\geometry{margin=2.5cm}

% Theorem environments (matching papers/su3_triplet/arxiv_version.tex)
\newtheorem{theorem}{Theorem}[section]
\newtheorem{lemma}[theorem]{Lemma}
\newtheorem{proposition}[theorem]{Proposition}
\newtheorem{corollary}[theorem]{Corollary}
\newtheorem{definition}[theorem]{Definition}
\theoremstyle{remark}
\newtheorem{remark}[theorem]{Remark}

% Notation shorthands
\newcommand{\B}{\mathcal{B}}
\newcommand{\C}{\mathbb{C}}
\newcommand{\R}{\mathbb{R}}
\newcommand{\H}{\mathbb{H}}
\newcommand{\Z}{\mathbb{Z}}
\newcommand{\Vc}{V_c}

\title{%
  Obstruction to Three Isomorphic Subspaces in $\mathbb{C}\otimes\mathbb{H}$\\
  \large Three Generations Research Track --- ST-2
}

\author{David Jaro\v{s}}

\date{2026-03-03}

\begin{document}

\maketitle

\begin{abstract}
We prove that the biquaternion algebra $\mathcal{B} = \mathbb{C}\otimes_\mathbb{R}\mathbb{H}
\cong M_2(\mathbb{C})$ does not admit a decomposition into three mutually
isomorphic subspaces that are simultaneously (a)~preserved by the $\mathbb{Z}_2^3$
grading of $\mathcal{B}$ and (b)~carry independent faithful representations of
$\mathrm{SU}(3)_{V_c}$.  Two independent obstructions are established: a
dimension obstruction for minimal left ideals (Theorem~\ref{thm:ideal_obstruction})
and a group-order obstruction for $\mathbb{Z}_2^3$-compatible decompositions
(Theorem~\ref{thm:grading_obstruction}).  Together these show that three
fermionic generations cannot arise from any purely algebraic decomposition of
$\mathbb{C}\otimes\mathbb{H}$ alone.  A UBT-specific mechanism via complex time
(AXIOM~B) is not precluded and is studied separately in ST-3
(\texttt{st3\_complex\_time\_generations.tex}).
\end{abstract}

\tableofcontents

%=============================================================
\section{Background and Statement of the Problem}
\label{sec:background}
%=============================================================

Furey~\cite{Furey2016} derives three generations of fermions using the
octonion-based algebra $\mathbb{C}\otimes\mathbb{O}$, exploiting its
decomposition into three minimal left ideals each isomorphic to $\mathbb{C}^4$.
The question addressed here is:

\begin{quote}
\emph{Does $\mathcal{B} = \mathbb{C}\otimes_\mathbb{R}\mathbb{H} \cong
M_2(\mathbb{C})$ admit any decomposition into three mutually isomorphic
subspaces that could serve as the algebraic seat of three fermion generations?}
\end{quote}

We give a definitive negative answer for all decompositions that must be
compatible with the established structure of~$\mathcal{B}$
(Theorems~\ref{thm:ideal_obstruction} and~\ref{thm:grading_obstruction}).

Throughout this document we use the notation and results of the companion
paper~\cite{su3triplet}: $\mathcal{B}=\mathbb{C}\otimes_\mathbb{R}\mathbb{H}$,
$\mathcal{B}_{\mathrm{basis}} = \{1,\mathbf{I},\mathbf{J},\mathbf{K},
i,i\mathbf{I},i\mathbf{J},i\mathbf{K}\}$, and $V_c =
\mathrm{span}_\mathbb{C}\{\mathbf{I},\mathbf{J},\mathbf{K}\}$ is the canonical
carrier space of $\mathrm{SU}(3)_{V_c}$.

%=============================================================
\section{Ideal Decomposition Obstruction}
\label{sec:ideals}
%=============================================================

\subsection{Minimal Left Ideals of $M_2(\mathbb{C})$}

\begin{definition}[Minimal left ideal]
  A left ideal $L \subset M_2(\mathbb{C})$ is \emph{minimal} if it is
  non-zero and contains no proper non-zero sub-ideal.
\end{definition}

\begin{lemma}[Structure of minimal left ideals]
  \label{lem:mli_structure}
  Every minimal left ideal of $M_2(\mathbb{C})$ is of the form
  $L_e = M_2(\mathbb{C})\cdot e$ for a primitive idempotent $e$
  (i.e.\ $e^2 = e$, $\mathrm{rank}(e) = 1$).
  Any two such ideals are isomorphic as left $M_2(\mathbb{C})$-modules,
  and each has $\dim_\mathbb{C} L_e = 2$.
\end{lemma}

\begin{proof}
  Standard Wedderburn--Artin theory: $M_2(\mathbb{C})$ is a simple
  $\mathbb{C}$-algebra of rank $2$ over $\mathbb{C}$, with unique simple
  left module $\mathbb{C}^2$.  Minimal left ideals correspond bijectively to
  primitive idempotents up to conjugacy; each is isomorphic to $\mathbb{C}^2$.
  The dimension count gives $\dim_\mathbb{C} L_e = 2$.
\end{proof}

\begin{theorem}[Ideal decomposition obstruction]
  \label{thm:ideal_obstruction}
  $M_2(\mathbb{C})$ cannot be written as a direct sum of three or more
  minimal left ideals.
\end{theorem}

\begin{proof}
  By Lemma~\ref{lem:mli_structure}, each minimal left ideal has
  $\dim_\mathbb{C} = 2$.  If $M_2(\mathbb{C}) = L_1 \oplus \cdots \oplus L_k$
  with each $L_j$ a minimal left ideal, then $\dim_\mathbb{C} M_2(\mathbb{C})
  = 2k$.  Since $\dim_\mathbb{C} M_2(\mathbb{C}) = 4$, we get $k = 2$.
  The standard decomposition $M_2(\mathbb{C}) = L_{e_1} \oplus L_{e_2}$
  with $e_1 = \mathrm{diag}(1,0)$, $e_2 = \mathrm{diag}(0,1)$ shows $k=2$
  is achieved.  Three ideals would require $\dim_\mathbb{C} \geq 6 > 4$,
  which is impossible.
\end{proof}

\begin{remark}
  The two minimal left ideals of $M_2(\mathbb{C})$ correspond to the two
  ``colour columns'' of the matrix algebra.  Each is 4-dimensional over $\mathbb{R}$.
  Any basis for three generations would need $3 \times 4 = 12$ real dimensions,
  which exceeds $\dim_\mathbb{R} \mathcal{B} = 8$.
\end{remark}

%=============================================================
\section{Grading-Compatible Decomposition Obstruction}
\label{sec:grading}
%=============================================================

\subsection{$\mathbb{Z}_2^3$ Eigenspace Structure}

Following~\cite{su3triplet}, the three commuting involutions $P_1, P_2, P_3$
on $\mathcal{B}$ produce six non-empty simultaneous eigenspaces:

\begin{center}
\begin{tabular}{ccccc}
  \hline
  Sector $(P_1,P_2,P_3)$ & Basis & $\dim_\mathbb{R}$ \\
  \hline
  $(+{,}+{,}+)$ & $\{1\}$ & $1$ \\
  $(-{,}+{,}+)$ & $\{i\}$ & $1$ \\
  $(+{,}-{,}+)$ & $\{\mathbf{J},\mathbf{K}\}$ & $2$ \\
  $(-{,}-{,}+)$ & $\{i\mathbf{J},i\mathbf{K}\}$ & $2$ \\
  $(+{,}+{,}-)$ & $\{\mathbf{I}\}$ & $1$ \\
  $(-{,}+{,}-)$ & $\{i\mathbf{I}\}$ & $1$ \\
  \hline
\end{tabular}
\end{center}

The total real dimension is $1+1+2+2+1+1 = 8 = \dim_\mathbb{R}\mathcal{B}$.
The carrier space of $\mathrm{SU}(3)_{V_c}$ is $V_c = \mathrm{span}_\mathbb{C}
\{\mathbf{I},\mathbf{J},\mathbf{K}\}$, which is the $P_2 = -1$ eigenspace
with $\dim_\mathbb{C} V_c = 3$.

\subsection{Obstruction from the Order of $\mathbb{Z}_2^3$}

\begin{lemma}[No $\mathbb{Z}_3$ in $\mathbb{Z}_2^3$]
  \label{lem:no_z3}
  The group $\mathbb{Z}_2^3$ contains no subgroup of order $3$.
\end{lemma}

\begin{proof}
  Every element of $\mathbb{Z}_2^3$ has order $1$ or $2$;
  $|\mathbb{Z}_2^3| = 8$ is not divisible by $3$.  By Lagrange's theorem,
  $\mathbb{Z}_2^3$ has no element or subgroup of order $3$.
\end{proof}

\begin{theorem}[$\mathbb{Z}_2^3$-grading obstruction]
  \label{thm:grading_obstruction}
  There is no decomposition $\mathcal{B} = W_1 \oplus W_2 \oplus W_3$ of
  $\mathcal{B}$ (as $\mathbb{R}$-vector spaces) into three subspaces $W_j$
  satisfying:
  \begin{enumerate}[label=(\alph*)]
    \item Each $W_j$ is a non-zero direct sum of $\mathbb{Z}_2^3$-eigenspaces.
    \item The three pieces $W_1, W_2, W_3$ are mutually isomorphic as
          $\mathbb{R}$-vector spaces (hence must have equal dimension).
    \item Each $W_j$ carries a faithful action of $\mathrm{SU}(3)_{V_c}$
          that factors through the standard representation on $V_c$.
  \end{enumerate}
\end{theorem}

\begin{proof}
  We argue in three steps.

  \medskip\noindent\textbf{Step 1 (Dimension constraint).}
  Condition~(b) requires $\dim_\mathbb{R} W_j = d$ for all $j$ and
  $3d = \dim_\mathbb{R}\mathcal{B} = 8$.  But $3 \nmid 8$, so no integer $d$
  solves $3d = 8$.  Hence condition~(b) alone is incompatible with
  a partition of $\mathcal{B}$ into three equal-dimensional parts.

  \medskip\noindent\textbf{Step 2 ($\mathrm{SU}(3)_{V_c}$ constraint).}
  The standard (fundamental) representation of $\mathrm{SU}(3)$ on $\mathbb{C}^3$
  has real dimension $6$.  If each $W_j$ is to carry a faithful copy of this
  representation, then $\dim_\mathbb{R} W_j \geq 6$.  Three such pieces would
  require $\dim_\mathbb{R}\mathcal{B} \geq 18 \gg 8$.  This is impossible.

  \medskip\noindent\textbf{Step 3 (Grading compatibility).}
  By condition~(a), each $W_j$ is a union of $\mathbb{Z}_2^3$-eigenspaces.
  By Lemma~\ref{lem:no_z3}, there is no $\mathbb{Z}_3$-action on the
  eigenspace labels compatible with the $\mathbb{Z}_2^3$ grading.  Therefore,
  any ``rotation'' among three isomorphic pieces cannot be realised as a
  $\mathbb{Z}_2^3$-graded automorphism of $\mathcal{B}$.

  Steps~1--3 together show no decomposition satisfying (a), (b), (c)
  can exist.
\end{proof}

\begin{corollary}[No three-generation mechanism in $\mathcal{B}$ alone]
  \label{cor:no_three_gen}
  Any three-generation mechanism in the UBT framework must involve structure
  beyond the pure algebra $\mathcal{B} = \mathbb{C}\otimes\mathbb{H}$.
  In particular, it cannot arise from left ideal decomposition, from the
  $\mathbb{Z}_2^3$ grading, from $\mathrm{Aut}_\mathbb{R}(\mathcal{B})$,
  or from any finite-group action intrinsic to $\mathcal{B}$.
\end{corollary}

%=============================================================
\section{Contrast with Furey's Octonionic Construction}
\label{sec:furey_contrast}
%=============================================================

For comparison, we note why the same argument does \emph{not} obstruct
three generations in Furey's $\mathbb{C}\otimes\mathbb{O}$:

\begin{itemize}
  \item $\dim_\mathbb{R}(\mathbb{C}\otimes\mathbb{O}) = 16$;
    $16 / 3$ is not an integer either, so the dimension argument in
    Step~1 above still applies in principle.
  \item However, Furey works with the \emph{left action algebra}
    $\mathbb{C}\ell(\mathbb{C}\otimes\mathbb{O})$, which is not the algebra
    itself.  The three minimal left ideals arise in the action algebra, which
    has higher dimension and different module structure.
  \item $\mathbb{C}\otimes\mathbb{O}$ is non-associative.  Non-associativity
    allows module structures impossible in the associative case.
\end{itemize}

These points confirm that the obstructions proved here are genuine features of
the \emph{associative} algebra $\mathcal{B}$, not universal statements about
all division-algebra-based constructions.

%=============================================================
\section{Conclusion}
\label{sec:conclusion}
%=============================================================

\begin{theorem}[Main obstruction theorem]
  \label{thm:main}
  $\mathcal{B} = \mathbb{C}\otimes_\mathbb{R}\mathbb{H} \cong M_2(\mathbb{C})$
  does not admit a decomposition into three mutually isomorphic subspaces
  that are (a)~preserved by the $\mathbb{Z}_2^3$ grading and (b)~carry
  independent $\mathrm{SU}(3)_{V_c}$ representations.
\end{theorem}

\begin{proof}
  Immediate from Theorem~\ref{thm:ideal_obstruction} (ideal obstruction)
  and Theorem~\ref{thm:grading_obstruction} (grading obstruction).
\end{proof}

\begin{remark}[What is not obstructed]
  Theorem~\ref{thm:main} does not obstruct three-generation mechanisms that
  involve structure \emph{external} to the pure algebra $\mathcal{B}$,
  such as the imaginary-time sector $\psi$ of AXIOM~B
  ($\tau = t + i\psi$).  The $\psi$-mode mechanism is studied in
  ST-3 (\texttt{st3\_complex\_time\_generations.tex}) and is not
  subject to the obstructions proved here.
\end{remark}

\begin{thebibliography}{99}

\bibitem{su3triplet}
D.~Jaro\v{s},
\emph{A Natural $\mathrm{SU}(3)$ Action on a Canonical Subspace of
$\mathbb{C}\otimes\mathbb{H}$ via $\mathbb{Z}_2^3$ Involutions},
\texttt{papers/su3\_triplet/arxiv\_version.tex}, 2026.

\bibitem{Furey2016}
C.~Furey,
Standard Model Physics from an Algebra?,
PhD Thesis, U.~Waterloo, 2015; \texttt{arXiv:1611.09182}.

\bibitem{WedArtin}
E.~Artin,
\emph{Geometric Algebra}, Interscience, 1957.

\bibitem{Porteous1995}
I.~R. Porteous,
\emph{Clifford Algebras and the Classical Groups}, Cambridge U.P., 1995.

\end{thebibliography}

\end{document}
